% Compile with: uplatex file.tex && dvipdfmx file.dvi
%           or: latexmk -uplatex file.tex
\documentclass[dvipdfmx]{article}
\usepackage{amsmath, amssymb}
\usepackage{tikz}

\title{scenario\_ito\_black\_scholes}
\date{\today}

\begin{document}
\maketitle

Apply Itô's lemma to a function of an asset following geometric Brownian motion $dS=\mu S\,dt+\sigma S\,dW$. Start by registering the logarithm $f=\ln S$.

$f = \ln S$

\begin{equation}
  f = \ln S
  \label{eq:E1}
\end{equation}

Obtain, with general operations, the three partial derivatives Itô's formula needs. First $\partial f/\partial t$, which is 0 since $t$ does not appear explicitly in $f$. Whether this is a partial or a total derivative is stated by the researcher (the ambiguity is real, since $S$ depends on $t$).

take the partial derivative of E1 with respect to t

\begin{equation}
  \frac{\partial f}{\partial t} = 0
  \label{eq:E2}
\end{equation}

Then $\partial f/\partial S$.

take the partial derivative of E1 with respect to S

\begin{equation}
  \frac{\partial f}{\partial S} = \frac{1}{S}
  \label{eq:E3}
\end{equation}

The second-order term $\partial^2 f/\partial S^2$ specific to Itô's lemma — this is what separates it from the classical chain rule.

take the 2nd order partial derivative of E1 with respect to S

\begin{equation}
  \frac{\partial^{2} f}{\partial S^{2}} = -\frac{1}{S^2}
  \label{eq:E4}
\end{equation}

Assemble Itô's formula $df=(\partial_t f+\mu S\partial_S f+\frac12\sigma^2S^2\partial_S^2 f)dt+\sigma S\partial_S f\,dW$ from those expression references. The formula itself is finance, and the researcher writes it.

$df = (E2 + \mu S E3 + \frac{1}{2}\sigma^2 S^2 E4) dt + \sigma S E3 dW$

\begin{equation}
  df=dt \left(\mu-\frac{\sigma^2}{2}\right)+dW \sigma
  \label{eq:E5}
\end{equation}

Repeat the same steps for a different function $g=S^2$, to confirm the same formula works on another input.

$g = S^2$

\begin{equation}
  g = S^2
  \label{eq:E6}
\end{equation}

$\partial g/\partial t$.

take the partial derivative of E6 with respect to t

\begin{equation}
  \frac{\partial g}{\partial t} = 0
  \label{eq:E7}
\end{equation}

$\partial g/\partial S$.

take the partial derivative of E6 with respect to S

\begin{equation}
  \frac{\partial g}{\partial S} = 2 S
  \label{eq:E8}
\end{equation}

$\partial^2 g/\partial S^2$, a constant 2.

take the 2nd order partial derivative of E6 with respect to S

\begin{equation}
  \frac{\partial^{2} g}{\partial S^{2}} = 2
  \label{eq:E9}
\end{equation}

Assemble the same Itô formula for $g$ — identical in shape to step 5, only the expression references change.

$dg = (E7 + \mu S E8 + \frac{1}{2}\sigma^2 S^2 E9) dt + \sigma S E8 dW$

\begin{equation}
  dg=dt \left(S^2 \sigma^2+2 S^2 \mu\right)+2 S^2 dW \sigma
  \label{eq:E10}
\end{equation}

The SDE notation itself needs no dedicated operation — it registers as an ordinary expression ($dt$ and $dW$ are plain symbols).

$dS = \mu S dt + \sigma S dW$

\begin{equation}
  dS = \mu S dt + \sigma S dW
  \label{eq:E11}
\end{equation}

The Black–Scholes equation is likewise an ordinary equation. Solving it symbolically for a closed-form price is out of this scenario's scope.

$\frac{\partial V}{\partial t} + \frac{1}{2}\sigma^2 S^2 \frac{\partial^2 V}{\partial S^2} + r S \frac{\partial V}{\partial S} - r V = 0$

\begin{equation}
  \frac{\partial V}{\partial t} + \frac{1}{2}\sigma^2 S^2 \frac{\partial^2 V}{\partial S^2} + r S \frac{\partial V}{\partial S} - r V = 0
  \label{eq:E12}
\end{equation}

\end{document}